\documentclass[11pt]{article}
\usepackage{amsmath,amssymb,amsthm}
\newtheorem{proposition}{Proposition}
\title{A Positive Perspective Reformulation of Admissible Scores}
\author{}
\date{November 24, 2026}
\begin{document}
\maketitle

\paragraph{Status.}
The following is an exact written reformulation, not a new upper or lower
bound for \texttt{worstCaseHeadComplexity}. It applies to every sum of
strictly oriented positive affine ratios, and therefore to the strict
ratio representations used by the unchanged model. It neither assumes
nor proves a small number of sign regions. The elementary identities
are checked on a rational example by
\texttt{positive-perspective-check.py}; the general proof is written,
not Lean-formalized.

Fix a head count $H\geq1$. For each head choose either $y_{hi}(x)=x_i$
for all $i$, or $y_{hi}(x)=1-x_i$ for all $i$. Write
\[
 B_h(x)=b_{h0}+\sum_{i=1}^n b_{hi}y_{hi}(x),\qquad
 A_h(x)=a_{h0}+\sum_{i=1}^n a_{hi}y_{hi}(x),\qquad
 b_{hi}>0\quad(0\leq i\leq n),
\]
and let $F(x)=c_0+\sum_h A_h(x)/B_h(x)$. This includes both allowed
strict denominator orientations. An everywhere-positive negative-slope
form is put in this form by complementing its coordinates; its new
constant is its positive value at the all-one vertex. The usual final
strict-slope perturbation, already checked in
\texttt{HeadComplexity.Results.AxisZeroGluing}, handles weak slopes when
only a strict finite sign pattern is required. A Boolean representation
with zero scores on false inputs is first made strict by subtracting a
sufficiently small positive bias, using finiteness of the cube. This note does not change
that existing admissibility argument.

\begin{proposition}[Positive coefficients and convex parameter sections]
There are real numbers $c_h$ and nonnegative numbers $u_{hi}$ such that
\[
 F(x)=\sum_{h=1}^H
 \frac{u_{h0}^2+\sum_i u_{hi}^2 y_{hi}(x)}{B_h(x)}-C,
 \qquad C=\sum_h c_h-c_0.
 \tag{1}
\]
For a fixed Boolean vertex $x$ and fixed choices of orientations, the
set of parameters $(u,b,C)$ with $b_{hi}>0$ and $F(x)<0$ is convex.
If $F$ is negative at any vertex, then $C>0$, and division by $C$
normalizes the threshold in (1) to one without changing any strict sign.
\end{proposition}
\begin{proof}
Choose
\[
 c_h\geq\max_{0\leq i\leq n}(-a_{hi}/b_{hi}),\qquad
 u_{hi}=\sqrt{a_{hi}+c_hb_{hi}}.
\]
Thus $A_h+c_hB_h=u_{h0}^2+\sum_i u_{hi}^2y_{hi}$, which proves (1).
At a fixed vertex put
\[
 q_h=(u_{h0},u_{h1}y_{h1}(x),\ldots,u_{hn}y_{hn}(x)).
\]
Because $y_{hi}(x)^2=y_{hi}(x)$, the summand in (1) is
$\|q_h\|^2/B_h(x)$. The quadratic perspective $(q,B)\mapsto\|q\|^2/B$
is convex on $B>0$: its second differential is
\[
 \frac{2}{B}\left\|dq-\frac{q}{B}\,dB\right\|^2\geq0.
 \tag{2}
\]
Here $q_h$ is linear in $u$ and $B_h(x)$ is linear in $b$.
Consequently $\sum_h\|q_h\|^2/B_h(x)-C$ is jointly convex, proving
the sublevel assertion. Every summand is nonnegative, so a negative
value of $F$ forces $C>0$. Replacing each $u_{hi}$ by $u_{hi}/\sqrt C$
gives $F/C$ with threshold one.
\end{proof}

\begin{proposition}[Convex extension and an affine matrix pencil]
The function
\[
 \widetilde F(x)=\sum_h
 \frac{u_{h0}^2+\sum_i u_{hi}^2y_{hi}(x)^2}{B_h(x)}-C
 \tag{3}
\]
is convex on $[0,1]^n$ and equals $F$ at every Boolean vertex.
Its strict negative sublevel is given by an explicit affine positive
matrix inequality. Namely, stack the vectors
$q_h(x)=(u_{h0},u_{h1}y_{h1}(x),\ldots,u_{hn}y_{hn}(x))$ into $q(x)$,
and put
\[
 D(x)=\operatorname{diag}(B_1(x)I_{n+1},\ldots,B_H(x)I_{n+1}),
 \qquad
 L(x)=\begin{pmatrix}D(x)&q(x)\\q(x)^{\mathsf T}&C\end{pmatrix}.
 \tag{4}
\]
Then $L(x)$ is affine in $x$ and
\[
             L(x)\succ0\quad\Longleftrightarrow\quad
             \widetilde F(x)<0.
 \tag{5}
\]
\end{proposition}
\begin{proof}
Each denominator stays positive on the entire continuous cube, and
both $q_h(x)$ and $B_h(x)$ are affine in $x$. Equation (2) proves
convexity. Equality at vertices is the Boolean identity $y^2=y$.
The matrix $D(x)$ is positive definite. Completing the square in the
quadratic form of $L(x)$ gives the Schur complement
$C-q(x)^{\mathsf T}D(x)^{-1}q(x)=-\widetilde F(x)$, proving (5).
All entries in (4) are affine in $x$.
\end{proof}

\paragraph{The exact Boolean-vanishing freedom.}
The continuous extension is not canonical. Increasing one shift $c_h$
by $t\geq0$, and adjusting the squared coefficients and $C$ accordingly,
changes it by exactly
\[
 -t\,\frac{\sum_i b_{hi}y_{hi}(x)(1-y_{hi}(x))}{B_h(x)}.
 \tag{6}
\]
Indeed the change in the numerator of (3) is
$t(b_{h0}+\sum_i b_{hi}y_{hi}^2)$ and the threshold increases by $t$.
The difference is (6). It vanishes at every vertex and is strictly
negative at every nonvertex of the cube, since all slopes are positive.
On any compact set disjoint from the vertices, its negative magnitude
has a positive minimum. Hence, for any fixed score, large shifts make
the convex extension negative throughout that compact set without
changing a single Boolean score. This freedom must be retained in any
geometric counting argument.

\paragraph{What this does not prove.}
An arbitrary collection of convex subsets of a parameter space can
have arbitrarily large dual shattering complexity; convexity alone
therefore does not give the desired $O(nH)$ sign-count exponent.
Likewise every subset of cube vertices has a convex separator, so the
existence of a convex extension is not a Boolean lower bound. The
matrix in (4) has special repeated scalar blocks. A useful theorem
would have to exploit this shared structure quantitatively.
Its determinant still contains the familiar cleared numerator:
\[
 \det L(x)=\left(\prod_h B_h(x)^{n+1}\right)
 \left(C-\sum_h\frac{\|q_h(x)\|^2}{B_h(x)}\right).
\]
Removing the positive factor $\prod_hB_h(x)^n$ recovers a polynomial
of degree at most $H+1$ in the ordinary affine numerator and denominator
coefficients. Thus merely introducing the matrix pencil does not remove
the degree factor in the existing Warren argument.

The original sum of heads and its head count are unchanged. No new
hypothesis establishing either global asymptotic target is supplied.
\end{document}
